\subsection{Mac Lane's separability criterion}

THEOREM 1. --- Let K be a field of characteristic 0. Every reduced K-algebra and in particular every extension of K is separable over K.

We first show that every extension L of K is separable. Let B be a transcendence basis of L over K (V, p.~109, Th. 1) and let \(L_{1} = \mathbf{K}(\mathbf{B})\) . Then \(L_{1}\) is separable over K (V, p.~121, Prop. 6). Further, L is an algebraic extension of \(L_{1}\) and \(L_{1}\) is a field of characteristic 0; therefore L is separable over \(L_{1}\) (V, p.~37, Cor.). By Prop. 9, L is thus separable over K.

Now let A be a reduced algebra over the field K. By the Cor. of Prop. 2 (V, p.~119) there exists a family \((L_{i})_{i \in I}\) of extensions of K such that A is isomorphic to a subalgebra of \(\prod_{i \in I} L_{i}\) . Each of the algebras \(L_{i}\) is separable over K by what has

been said and so A has the same property, by Prop. 3 a) and c) (V, p.~119).

THEOREM 2. --- Let K be a field of characteristic \(p \neq 0\) , \(\mathbf{K}^{p^{-\infty}}\) a perfect closure of K and A a commutative K-algebra. The following properties are equivalent:

\begin{enumerate}
\def\labelenumi{\alph{enumi})}
\item
  A is separable.
\item
  There exists an extension \(\mathbf{K}'\) of \(\mathbf{K}\) such that \(\mathbf{K}'\) is perfect and \(\mathbf{K}' \otimes_{\mathbf{K}} \mathbf{A}\) is reduced.
\item
  The ring \(\mathbf{K}^{p^{-\infty}}\otimes_{\mathbf{K}}\mathbf{A}\) is reduced.
\item
  The ring \(\mathbf{K}' \otimes_{\mathbf{K}} \mathbf{A}\) is reduced for every extension \(\mathbf{K}'\) of \(\mathbf{K}\) which is of finite degree and \(p\) -radical of height \(\leqslant 1\) .
\item
  For every family \((a_{i})_{i\in I}\) of elements of A linearly free over K, the family \((a_{i}^{p})_{i\in I}\) is linearly free over K.
\item
  There exists a basis \((a_{i})_{i\in I}\) of the vector K-space A such that the family \((a_{i}^{p})_{i\in I}\) is linearly free over K.
\end{enumerate}

If an extension \(\mathbf{K}'\) of \(\mathbf{K}\) is a perfect field, it contains a subextension \(\mathbf{K}\) -isomorphic to \(\mathbf{K}^{p^{-\infty}}\) (V, p.~6, Prop. 3); moreover, every \(p\) -radical extension of \(\mathbf{K}\) is isomorphic to a subextension of \(\mathbf{K}^{p^{-\infty}}\) (V, p.~26, Prop. 3). These remarks show the implications \(a) \Rightarrow b) \Rightarrow c) \Rightarrow d)\) .

Let us prove that \(d)\) implies \(e)\) . Let \((a_i)_{i \in I}\) be a linearly free family in \(A\) and let \((\lambda_i)_{i \in I}\) be a family of finite support in \(K\) such that \(\sum \lambda_i a_i^p = 0\) . Let \(K'\) be the subextension of \(K^{p^{-\infty}}\) generated by the elements \(\lambda_i^{p^{-1}}\) ; it is of finite degree and height \(\leqslant 1\) . Put \(x = \sum_{i \in I} \lambda_i^{p^{-1}} \otimes a_i\) in \(K' \otimes_K A\) ; we have

\[
x ^ {p} = \sum_ {i \in I} \lambda_ {i} \otimes a _ {i} ^ {p} = 1 \otimes \sum_ {i \in I} \lambda_ {i} a _ {i} ^ {p} = 0.
\]

By the hypothesis \(d)\) we have \(x = 0\) , whence \(\lambda_i = 0\) for all \(i \in I\) .

Clearly \(e)\) implies \(f)\) and it remains to show that \(f)\) implies \(a)\) . Thus let \((a_i)_{i \in I}\) be a basis of \(A\) over \(K\) such that the family \((a_i^p)_{i \in I}\) is linearly free over \(K\) . Let \(L\) be an extension of \(K\) and let \(x\) be an element of \(L \otimes_K A\) such that \(x^p = 0\) . Write \(x = \sum_{i \in I} \lambda_i \otimes a_i\) with \(\lambda_i \in L\) for all \(i \in I\) . We have \(x^p =\)

\(\sum_{i\in I}\lambda_i^p\otimes a_i^p = 0\) and since the family \((a_i^p)_{i\in I}\) is linearly free over \(\mathbf{K}\) , we have

\(\lambda_{i}^{p} = 0\) , whence \(\lambda_{i} = 0\) for all \(i \in I\) ; it follows that \(x = 0\) . So we have proved that \(x^{p} = 0\) implies \(x = 0\) in \(\mathbf{L} \otimes_{\mathbf{K}} \mathbf{A}\) , from which it follows at once that \(\mathbf{L} \otimes_{\mathbf{K}} \mathbf{A}\) is reduced.

COROLLARY 1 (Mac Lane). --- Let K be a field of characteristic exponent p, Ω a perfect extension of K and L a subextension of Ω. Then the following conditions are equivalent :

\begin{enumerate}
\def\labelenumi{\alph{enumi})}
\item
  L is separable over K.
\item
  L is linearly disjoint from \(\mathbf{K}^{p^{-\infty}}\) over \(\mathbf{K}\) .
\item
  L is linearly disjoint over K from every p-radical extension of K contained in \(\Omega\) , of finite degree and height \(\leqslant 1\) .
\end{enumerate}

The case where K is of characteristic 0 is trivial because then L is separable over K (Th. 1) and \(K^{p^{-\infty}} = K\) by convention. Suppose then that \(p \neq 1\) , and let us first show that a) implies b). Assume that L is separable over K and let \((a_i)_{i \in I}\) be a basis of L over K. Let \((\lambda_i)_{i \in I}\) be a family of finite support of elements of \(K^{p^{-\infty}}\) such that \(\sum_{i \in I} \lambda_i a_i = 0\) ; there exists an integer \(f \geqslant 0\) and elements

\(\mu_{i}\) of \(\mathbf{K}\) such that \(\lambda_{i} = \mu_{i}^{p^{-f}}\) . We have

\[
\sum_ {i \in I} \mu_ {i} a _ {i} ^ {p ^ {f}} = \left(\sum_ {i \in I} \lambda_ {i} a _ {i}\right) ^ {p ^ {f}} = 0
\]

and Th. 2 implies, by induction on \(f\) , that the family \((a_i^{p^f})_{i\in I}\) is linearly free over \(K\) . We thus have \(\mu_i = 0\) , whence \(\lambda_i = 0\) for all \(i\in I\) . Finally \(L\) is linearly disjoint from \(K^{p^{-\infty}}\) over \(K\) .

It is clear that \(b)\) implies \(c)\) . Lastly suppose that \(c)\) holds and let \(K'\) be an extension of \(K\) of finite degree and \(p\) -radical of height \(\leqslant 1\) . The ring \(K' \otimes_K L\) is isomorphic to a subring of \(\Omega\) , hence is reduced. Now it follows from Th. 2 that \(L\) is separable over \(K\) .

COROLLARY 2. --- Let K be a field of characteristic exponent p, \(K^{p^{-\infty}}\) a perfect closure of K and L a separable extension of K. Then the ring \(L \otimes_{K} K^{p^{-\infty}}\) is a field. If moreover L is algebraic over K, then \(L \otimes_{K} K^{p^{-\infty}}\) is a perfect closure of L.

The case \(p = 1\) is trivial, so let us suppose \(p \neq 1\) . Let \(\Omega\) be a perfect closure of \(L\) . By Cor. 1 there exists a K-algebra homomorphism of \(L \otimes_K K^{p^{-\infty}}\) onto \(L[K^{p^{-\infty}}]\) which maps \(x \otimes y\) to \(xy\) for \(x \in L\) and \(y \in K^{p^{-\infty}}\) . Since \(K^{p^{-\infty}}\) is algebraic over \(K\) , the ring \(L[K^{p^{-\infty}}]\) is a subfield of \(\Omega\) (V, p.~18, Cor. 1). Suppose further that \(L\) is algebraic over \(K\) , then \(L[K^{p^{-\infty}}]\) is an algebraic extension of the perfect field \(K^{p^{-\infty}}\) , hence it is a perfect field (V, p.~43, Cor. 1); finally, since the field \(L[K^{p^{-\infty}}]\) is a \(p\) -radical extension of \(L\) (V, p.~25, Cor.), it is a perfect closure of \(L\) .

COROLLARY 3.--- Let L be an algebraic extension of K and M a commutative L-algebra (for example an extension of L). If M is separable over K, it is separable over L.

The algebra L is K-isomorphic to a subalgebra of M, hence L is a separable extension of K. Therefore (Cor. 2) there exists an L-isomorphism of \(L^{p^{-\infty}}\) onto \(K^{p^{-\infty}} \otimes_{K} L\) . The ring \(L^{p^{-\infty}} \otimes_{L} M\) is thus isomorphic to \((\mathrm{K}^{p^{-\infty}} \otimes_{K} L) \otimes_{L} M\) , and so to \(K^{p^{-\infty}} \otimes_{K} M\) (II, p.~278, Prop. 2) and this latter ring is reduced because M is separable over K. Thus the ring \(L^{p^{-\infty}} \otimes_{L} M\) is reduced, which proves that M is separable over L (V, p.~122, Th. 2).

Remark. --- Mac Lane's criterion may be formulated without introducing any extensions of K other than L. For by c) of Cor. 1, L is separable over K if and only if L and \(K^{p^{-1}}\) are linearly disjoint over K. Since the mapping \(x \mapsto x^{p}\) is an isomorphism of L onto the subfield \(L^{p}\) , we obtain the following criterion (cf.~V, p.~177, Ex. 11 for an analogous criterion for algebras):

L is separable over K if and only if the subfields \(L^{p}\) and K of L are linearly disjoint over \(K^{p}\) .

